% Teil II: der durchgerechnete Weg. Reihenfolge beibehalten.

\subsection*{a) $f(x)=\frac{\sin x}{x}$ mit $f(0)$ passend}
Die Funktion ist für $x\neq0$ durch $\sin x/x$ gegeben; gesucht ist die
Fortsetzung an $x=0$, damit $f\in C^1(\mathbb R)$ liegt. Die Fortsetzung
$f(0)=1$ ergibt sich aus
\begin{equation}
  \lim_{x\to0}\frac{\sin x}{x}=1
\end{equation}
(bekannter Standardgrenzwert bzw.\ de l'Hôpital). Für $x\neq0$ ist mit der
Quotientenregel
\begin{equation}
  f'(x)=\frac{x\cos x-\sin x}{x^2}.
\end{equation}
An der Stelle $0$ prüfen wir die Ableitung über die Definition (bzw.\ als
Grenzwert von $f'$). Wegen $\sin x = x-\frac{x^3}{6}+\cdots$ und
$\cos x=1-\frac{x^2}{2}+\cdots$ ist
\begin{equation}
  x\cos x-\sin x
  = x\left(1-\frac{x^2}{2}\right)-\left(x-\frac{x^3}{6}\right)
  = -\frac{x^3}{3}+\frac{x^3}{6}
  = -\frac{x^3}{6}
  \;\text{(führender Term)}.
\end{equation}
Damit folgt für den Grenzwert der Ableitungsfunktion
\begin{equation}
  \lim_{x\to0} f'(x)
  =\lim_{x\to0}\frac{x\cos x-\sin x}{x^2}
  =\lim_{x\to0}\frac{-x^3/6}{x^2}
  =0.
\end{equation}
\begin{ergebnis}
\begin{equation*}
  f(0)=1,\qquad f'(0)=0,\qquad f\in C^1(\mathbb R).
\end{equation*}
\end{ergebnis}
Die konkrete Ableitung in $0$ über die Definition 4.1.1 (Differenzenquotient)
nachrechnen; das Ergebnis $0$ macht $f'$ stetig.

\subsection*{b) Fallunterscheidung an der Nahtstelle}
Die genaue Vorschrift der zweiten Funktion ist im eingescannten Bild unlesbar.
Methode: Für jede Seite die Ableitung über die Regeln bilden und an der
Nahtstelle beide einseitigen Ableitungswerte über die Definition vergleichen;
stimmen sie überein, ist $f'$ da (eindeutig) definiert. Zusätzlich die
Stetigkeit von $f'$ an der Nahtstelle prüfen. Bitte die konkreten Funktionsteile
am Blatt bestätigen und damit durchrechnen.

\begin{kontrolle}
\textbf{Querprobe:} $f'$ an der Nahtstelle einmal als Grenzwert der Formel,
einmal über den Differenzenquotienten — beide müssen denselben Wert liefern.
\textbf{Grössenordnung:} $\sin x/x\to1$ und $f'(0)\to0$ sind bekannte,
plausible Grenzwerte.
\end{kontrolle}

\begin{merksatz}
$C^1$ heisst: differenzierbar überall \emph{und} Ableitung stetig. An
Nahtstellen über die Definition der Ableitung prüfen; hebbare Lücken
(z. B. $\sin x/x$) durch den passenden Funktionswert schliessen.
\end{merksatz}
